% !TEX encoding = UTF-8 Unicode
% !TEX root = ../Σαούκιν.tex

\chapter{Introduction} \label{c1:intro}

Scientific research relies on previously published work of other researchers in the same or adjacent domains. A rare research does not build upon any work of the scientific community produced before it.

When publishing a scientific research paper (terms publication and [scientific] paper in this document refer to the same meaning) \snote{do not overload terms, use a single terminology}, researchers list references used to produce their work.
These references are called citations.

References help provide authority to the foundational ideas, identify intersections with relevant work of other researchers.
At the same time, unreferenced quotations can be considered plagiarism.

The extent to which the work of other authors has been cited in any given paper \scor{isn't}{is not} represented in the list of references.

The reverse - a count of all citations of a single author/paper in all other publications - is called the citations index.
A variety of formulas to calculate the citation index has been proposed.

Coincidentally, the citation index can provide an indication of significance of said author or paper.
While the specifics of different indices address concerns, the main idea is that the number of citations and the researcher impact are directly related and grow proportionally.

Authors can cite their own previously published work - this is called self-citation.
While self-citations aren't inherently bad (in the new or unpopular areas of research lacking active development, it's more probable to reuse some of the previously published own work), it can be used by bad actors (authors) to inflate their citation numbers and to provide authority to otherwise unsubstantial research areas or works.

A new metric proposed and introduced is the count of self-citations - the self-citation index.

\snote{I do not like the above. You mix way too many things. There is no need to. Keep things simple. For instance,
\begin{itemize}
\item
Research publications reference related publications. If a publication $x$ is referenced in a publication $y$ we say that $y$ \emph{is citing} $x$. 
\item
For each publication $x$, we are interested to know the number of its \emph{citations}, i.e., the number of publications $y$ that refer to $x$. 
\item
Citations may be distinguish into two categories: \emph{self} and \emph{third-party} citations. Define and explain.
\item
Typically, a publication with a higher number of third-party citations is more influential than a publication with a lower number of  third-party citations. Analogously, a research that has publications that have attracted a higher number of third-party citations is more influential than a researcher that has publications that attracted a lower number of third-party citations.
\item
Now, you may give a concrete example that measures the total citations of a researcher. See below:
\end{itemize}
}

\color{teal}

 Let us consider the researchers named Gabriel Ghinita, Panos Kalnis and Spiros Skiadopoulos that had published article [GPK2007] that is described as follows.
\begin{quote}
\item[{[GPK2007]}] Gabriel Ghinita, Panos Kalnis, Spiros Skiadopoulos:
\emph{PRIVE: anonymous location-based queries in distributed mobile systems}. WWW 2007: 371-380
\end{quote}
Panos Kalnis, a co-author of [GPK2007], later in the same year published article [KGMP2007] that cites [GPK2007] and is detailed as follows.
\begin{quote}
\item[{[KGMP2007]}] Panos Kalnis, Gabriel Ghinita, Kyriakos Mouratidis, Dimitris Papadias:
\emph{Preventing Location-Based Identity Inference in Anonymous Spatial Queries}. IEEE Trans. Knowl. Data Eng. 19(12): 1719-1733 (2007)
\end{quote}
It is clear that when we consider the third-party citations of Panos Kalnis and Gabriel Ghinita the citation of [GPK2007] in [KGMP2007] should be excluded (since Panos Kalnis and Gabriel Ghinita co-authored both articles).

For the case of Spiros Skiadopoulos the choice is not so clear. \snote{Carry on!}


\color{gray}

\snote{All gray text below is something that you should read carefully and respect (or better obey)!}

Η σημερινή εποχή χαρακτηρίζεται από ένα κόσμο ψηφιακό, όπου ένας τεράστιος αριθμός πληροφοριών, ...

Βασικά χαρακτηριστικά της εργασίας είναι:
\begin{itemize} 
\item
Η δυνατότητα άμεσης αλληλεπίδρασης με το χρήστη.
\item
...
\end{itemize}

Στις επόμενες σελίδες, παρουσιάζεται η ...

\section{Γενικά}
Το παρόν κεφάλαιο αποτελεί μια συνοπτική αναφορά ...

\section{Σκοπός}
Σκοπός αυτής της πτυχιακής είναι ...


\section{Δομή της εργασίας}

Η παρούσα εργασία δομείται ως εξής:
\begin{description}
\item[Στο πρώτο κεφάλαιο] της εργασίας, περιγράφονται ...

\item[Στο δεύτερο κεφάλαιο,] παρουσιάζονται  ...

\item[Στο τρίτο κεφάλαιο,] ...

\item[Στο ΧΧΧ και τελευταίο κεφάλαιο,] αναφέρονται μερικά συμπεράσματα σχετικά με την εργασία στο σύνολό της και ιδέες σχετικά με τις μελλοντικές επεκτάσεις της.
\end{description}





